\documentclass[11pt]{article}

\usepackage[margin=1in]{geometry}
\usepackage{amsmath,amssymb,amsthm,mathtools}
\usepackage{enumitem}
\usepackage{booktabs}
\usepackage{tabularx}
\usepackage{microtype}
\usepackage[hidelinks]{hyperref}
\usepackage{xurl}

\newtheorem{definition}{Definition}
\newtheorem{lemma}{Lemma}
\newtheorem{remark}{Remark}

\title{Challenge Stop Profiles for Successor Maintenance:\\[0.5ex]
\large Question Families, Semantic Pieces, and Approved Terminal Fields}
\author{}
\date{Unified archive note v0.08 (2026-03-28; archive v736)}

\begin{document}
\maketitle

\begin{abstract}
Audience-keyed challenge routes answer one practical question: where should a given reader start, and how should that route escalate?
They should not also own the semantic decomposition of each downstream question or the exact narrow fields where escalation is allowed to terminate.
This note isolates that missing layer.
It defines a question-keyed \emph{challenge stop profile} whose job is to decompose one recurring downstream question into contested semantic pieces and to name the approved terminal fields for each piece.
The result is a compact, auditable home for stop semantics that later papers can cite without repeating terminal-field lists inside audience ladders.
\end{abstract}

\paragraph{Non-redundancy note.}
Synthesis~31 owns the downstream normal form, Synthesis~34 owns the status-envelope / lineage-notice split, Synthesis~35 owns sentence locks, Synthesis~36 owns audience-keyed routes, Synthesis~38 owns piece-keyed challenge branches, Synthesis~54 owns the answer-side terminal citation normal form, Synthesis~74 owns the paired terminal citation discipline, Synthesis~75 owns the paired cutover rule, and Synthesis~17 owns the worked instantiation. This note owns only question families, semantic pieces, and approved terminal fields. Later notes should reopen it only when the exact question family, semantic-piece split, or approved terminal-field set is itself live.

\paragraph{Scope.}
This is not a new verdict, route, or receipt object. It is a publication-interface note about where a challenge may stop once the contested semantic piece is known.

\paragraph{Imports.}
Minimal citation targets are imported from Synthesis~31.\cite{series:downstreamnormalforms}
Audience-keyed start paths and escalation order are imported from Synthesis~36.\cite{series:challengeroutes}
The concrete worked instantiation is imported from Synthesis~17.\cite{series:workedexample}

\paragraph{Exports.}
This note exports (i) question-keyed challenge stop profiles, (ii) semantic-piece decompositions for recurring downstream questions, (iii) an approved-terminal-field rule preventing wrapper prose from quietly becoming a stopping point, and (iv) a local-reopen rule saying that later terse papers should reopen this note only when the exact question family, semantic-piece split, or approved terminal-field set is itself live. Otherwise later terse papers should stop at the answer-side terminal citation normal form in Synthesis~54, then stop at Synthesis~73 for the maintained request-side terminal stop, drop to Synthesis~17 for one concrete checked-answer/request-tail seam, and widen to Synthesis~74--75 only for the abstract paired terminal discipline or abstract stop/widen/reopen judgment. Exact piece-keyed escalation walks remain outside scope here and belong to Synthesis~38.\cite{series:challengebranches,series:challengeanswerreviewmenus,series:publicrequestresponsepacketrefreshresponsepacketclosureverdicts,series:pairedterminalcitationdiscipline,series:pairedterminalcutoverrules,series:workedexample}

\paragraph{One branched stop route.}
Within the same challenge-entry family, Synthesis~37 is the semantic-stop branch above one audience route: it says which contested pieces exist and which terminal fields may discharge them, but it does not choose a reader entry point or spell the exact contested-piece walk. Synthesis~35 remains the optional sentence-lock sidecar, Synthesis~36 remains the audience-entry owner, and Synthesis~38 remains the exact route-plus-piece owner. Stop here only while the live issue is which semantic piece or approved stop field is honest. Hand off left to Synthesis~36 when the audience entry itself moves, and right to Synthesis~38 when the exact branch path itself becomes live.\cite{series:challengeroutes,series:challengebranches}

\paragraph{One tiny stop-branch drill.}
For the worked compact-diff challenge, Synthesis~36 says where an auditor should begin, Synthesis~37 says that the recompute piece may stop at the verifier recomputation slice, and Synthesis~38 says what exact path reaches that slice. If a later terse paper only needs the approved stop field family, stopping at Synthesis~37 is honest; widening to the branch map before the path matters is unnecessary.\cite{series:challengeroutes,series:challengebranches,series:workedexample}

\section{Why the stop layer should be separate}
A route telling a referee to start at the closure ledger and an auditor to start at the status envelope does not yet answer a different question: which exact narrow fields are allowed to discharge the challenged semantic pieces of those questions?
If audience routes carry that terminal-field material inline, the same stop lists get repeated across multiple audiences and gradually drift.
So the semantic stopping layer deserves one small home of its own.

\begin{definition}[Challenge stop profile]
For one concrete base$\to$successor comparison, a \emph{challenge stop profile}
\[
\Sigma(q)=(P_q,F_q)
\]
for downstream question family $q$ consists of a finite nonempty semantic-piece set
\[
P_q=\{p_1,\dots,p_k\}
\]
and a finite nonempty approved-terminal-field family
\[
F_q=\{(p_i,f_{i1}),\dots,(p_i,f_{ir_i})\}_{i=1}^k,
\]
where each $f_{ij}$ is an already-owned narrow field or row slice whose semantics are sufficient to answer piece $p_i$.
\end{definition}

\begin{definition}[Approved-terminal-field rule]
A stop profile satisfies the \emph{approved-terminal-field rule} if every approved field in $F_q$ is narrower than the wrapper prose for $q$ and is already owned by one of the compare, continuity, closure, envelope, or ledger objects for the same successor comparison.
\end{definition}

\begin{definition}[Route-link rule]
An audience-keyed challenge route for question family $q$ satisfies the \emph{route-link rule} if it names the stop profile $\Sigma(q)$ rather than redefining approved terminal fields inline.
\end{definition}

\begin{definition}[First-sufficient-stop rule]
A stop profile satisfies the \emph{first-sufficient-stop rule} if, once a contested semantic piece $p_i$ is fixed, escalation may stop at the first approved field for $p_i$ whose owner directly discharges that piece; widening past that field is permitted only when a different semantic piece or request class is actually being asked.
\end{definition}

\begin{lemma}[Stop-profile reuse is audit-safe]
Assume a stop profile for question family $q$ satisfies the approved-terminal-field rule and the first-sufficient-stop rule, and every audience route for $q$ satisfies the route-link rule.
Then multiple audience routes may share the same terminal semantics without risking route-local drift about where a challenge may stop.
\end{lemma}

\begin{proof}
The approved-terminal-field rule forces the stopping points into already-owned narrow fields rather than wrapper prose.
The first-sufficient-stop rule says a challenge need not wander farther once the owner of the contested semantic piece has been reached.
The route-link rule then keeps those stopping points in one shared object instead of repeating them across audience ladders.
So different audiences may navigate differently while still inheriting the same terminal semantics.
\end{proof}

\begin{table}[t]
\centering
\small
\begin{tabularx}{\linewidth}{@{}l X X@{}}
\toprule
Question family & Semantic pieces & Typical approved terminal fields \\
\midrule
Public status sentence & continuity piece; closure piece & continuity verdict; closure verdict \\
Package completeness & role-row piece; closure token piece & closure ledger rows; closure-status token \\
Public status token & token piece; continuity piece; closure piece & status envelope token; continuity verdict; closure verdict \\
Compact diff summary & delta piece; recompute piece; changed-family piece & compare-report diff; verifier recompute slice; changed-line-items slice \\
\bottomrule
\end{tabularx}
\caption{Stop profiles are question-keyed semantic decompositions. Audience routes may start differently, but they should share these approved terminal fields.}
\label{tab:stopprofiles}
\end{table}

\section{Worked example pointer}
The maintained worked bundle now ships \path{example_successor_challenge_stop_profiles.json}.\cite{series:workedexample}
It records four question families for the canned Case-B successor: public status sentence, package completeness, public status token, and compact diff summary.
Each family lists its contested semantic pieces and the exact approved terminal fields that may answer them.
The companion route object then points at these stop profiles rather than repeating terminal-field lists in every audience ladder, and the public first-stop catalog now imports the same \texttt{stop\_profile\_key} for the four shipped audience/question answers so a later terse paper can reach the correct stop semantics directly from the maintained start object. The companion branch-map object then uses those same piece keys and terminal fields to record the exact piece-specific escalation walk for each audience route.\cite{series:challengeroutes,series:challengebranches}

\paragraph{Why this helps the archive.}
The series can now keep the citation map about smallest sufficient citation targets, the route note about where different readers should begin, this note about where specific contested semantic pieces may stop, and the new branch-map note about the exact piece-specific escalation walk.
That keeps public wrappers terse without letting their stop semantics go implicit.

\noindent That semantic stop layer becomes easier to reuse safely once the note says, in one place, which imported owner must be reopened first after a stale question-family shortcut fails.

\paragraph{A compact challenge-stop-profile first-reopen-owner card.}
\begin{table}[t]
\centering
\small
\begin{tabularx}{\linewidth}{@{}lX@{}}
\toprule
Field & First owner to reopen and why \\
\midrule
Checked-sentence reuse or inline-provenance drift & Reopen Synthesis~35 first. If the maintained shortcut no longer serves the same checked sentence or no longer inherits the same allowed unchanged-travel / recheck posture, the sentence-lock sidecar regains authority before this note can keep the same semantic stop claim honest. \\
Minimal citation-target drift & Reopen Synthesis~31 first. If the maintained shortcut no longer points at the same smallest sufficient downstream citation target, the downstream normal-form owner moves first before this note can keep the same semantic stop citation honest. \\
Audience-route drift & Reopen Synthesis~36 first. If the maintained shortcut no longer begins at the same audience-keyed start object or follows the same coarse escalation ladder, the route owner moves first before any approved stop profile can remain valid. \\
Question-family / semantic-piece / approved-terminal-field drift with imports fixed & Reopen Synthesis~37 first. If the same checked sentence, minimal citation target, and audience route stay fixed but the maintained question family changes, a semantic piece is added / removed / renamed, or an approved stop field changes, this note is the first honest owner again. \\
Exact route-plus-piece walk becomes live & Widen to Synthesis~38. Once the maintained semantic split and approved stop fields are fixed and the live issue becomes the exact audience-route-plus-piece walk to one approved field, the exact branch owner becomes the honest stop. \\
Coverage, exact surfaced realization, or exact narrow owner becomes live & Widen to Synthesis~39--41. Once the maintained stop profile is fixed and the live issue becomes family-level coverage, exact surfaced realization, or exact narrow terminal ownership, the later coverage / realization / terminal-owner notes regain authority. \\
Later answer-side stop or paired cutover becomes live & Stop at Synthesis~54, drop to Synthesis~17 for one concrete checked-answer/request-tail seam, or widen to Synthesis~74--75 only for the abstract paired terminal discipline or abstract stop/widen/reopen judgment. Once the maintained semantic split and approved stop fields are fixed, later terse papers should move there rather than reopen this note merely to restate a stable stop profile. \\
\bottomrule
\end{tabularx}
\caption{A compact challenge-stop-profile first-reopen-owner card. This is the smallest owner-cutover map later terse papers can quote directly when a maintained semantic stop-profile shortcut has gone stale and the remaining task is to recover which narrower or wider owner regains authority first.}
\label{tab:challenge-stop-profile-first-reopen-owner-card}
\end{table}

This card is intentionally fail-closed: it keeps Synthesis~37 narrow by routing checked-sentence and minimal-citation drift back to Synthesis~35 and Synthesis~31, audience-entry drift back to Synthesis~36, semantic-piece / approved-stop drift back to itself, and only then widening to the exact branch, coverage, realization, terminal-owner, or later answer-side cutover layers once the semantic stop profile is no longer the honest moving part.

\paragraph{A tiny first-reopen-owner drill.}
For the worked compact-diff question, the maintained stop profile remains an honest shortcut only while the same checked sentence may still travel under the same sentence-lock posture, the same smallest sufficient downstream citation target stays in view, the same audience route still imports the same stop-profile key, and the same semantic pieces still stop at the same approved fields \nolinkurl{observed_diffs}, \nolinkurl{recomputed}, and \nolinkurl{changed_line_items}. Reopen Synthesis~35 first if the compact-diff sentence itself is no longer licensed for the same reuse posture. Reopen Synthesis~31 first if the smallest sufficient citation target changes, and Synthesis~36 first if the audience-entry route changes. Reopen Synthesis~37 first if those imports stay fixed but the maintained compact-diff semantic split or approved stop fields change. Only after the semantic stop profile is fixed should a later terse paper widen to Synthesis~38 for the exact route-plus-piece walk, to Synthesis~39--41 for coverage / realization / terminal ownership, stop at Synthesis~54 for the maintained checked-answer handoff, then stop at Synthesis~73 for the maintained request-side terminal stop, drop to Synthesis~17 for one concrete checked-answer/request-tail seam, or widen to Synthesis~74--75 only for later abstract paired-terminal questions.\cite{series:workedexample,series:sentencelocks,series:downstreamnormalforms,series:challengeroutes,series:challengebranches,series:challengecoverage,series:challengesurfacehooks,series:challengeterminalwitnesses,series:challengeanswerreviewmenus,series:publicrequestresponsepacketrefreshresponsepacketclosureverdicts,series:pairedterminalcitationdiscipline,series:pairedterminalcutoverrules,series:workedexample}

\paragraph{Why this helps the archive.}
The series can now keep minimal citation targets, audience ladders, semantic stop profiles, and exact piece-local branches separate while also stating, in one narrow fail-closed card, which imported owner must be reopened first once a maintained semantic-stop shortcut goes stale.
That makes the worked example more auditable without lengthening the public wrapper sentence or forcing every downstream note to improvise its own semantic stop cutover from prose.


\paragraph{A minimal public challenge-stop-profile card.}
\begin{table}[t]
\centering
\small
\begin{tabularx}{\linewidth}{@{}lX@{}}
\toprule
Field & Minimal public answer \\
\midrule
Public-status sentence profile & \nolinkurl{public_status_sentence} splits into \nolinkurl{continuity_piece} and \nolinkurl{closure_piece}, which may stop only at \nolinkurl{continuity_decision} in \nolinkurl{example_successor_continuity_verdict.json} and \nolinkurl{closure_status} in \nolinkurl{example_publication_closure_verdict.json} \\
Package-completeness profile & \nolinkurl{package_completeness} splits into \nolinkurl{role_rows_piece} and \nolinkurl{closure_token_piece}, which may stop only at \nolinkurl{role_rows} in \nolinkurl{example_release_closure_ledger.json} and \nolinkurl{closure_status} in \nolinkurl{example_publication_closure_verdict.json} \\
Compact-diff profile & \nolinkurl{compact_diff_summary} splits into \nolinkurl{diff_delta_piece}, \nolinkurl{recompute_piece}, and \nolinkurl{changed_family_piece}, which may stop only at \nolinkurl{observed_diffs} in \nolinkurl{example_compare_report.json}, \nolinkurl{recomputed} in \nolinkurl{example_verifier_report.json}, and \nolinkurl{changed_line_items} in \nolinkurl{example_successor_continuity_verdict.json} \\
Selection rule & Stop here only while the live issue is which semantic piece belongs to one maintained question family and which approved narrow field may discharge that piece; do not widen to exact audience routing or exact contested-piece walks unless those neighboring owners have become the honest moving part \\
Left/right boundary & Reopen Synthesis~31 or Synthesis~36 when the live issue is the minimal citation target or audience start object; widen to Synthesis~38 once the exact piece-keyed escalation walk is the honest moving part, and to Synthesis~39--41 once the live issue has become coverage, exact surface realization, or exact terminal ownership rather than the stop profile itself \\
Paired-terminal boundary & Once the maintained semantic-piece split and approved terminal fields are fixed, later terse papers should usually stop at Synthesis~54, then stop at Synthesis~73 for the maintained request-side terminal stop, drop to Synthesis~17 for one concrete checked-answer/request-tail seam, and widen to Synthesis~74--75 only for the abstract paired-tail discipline or abstract stop/widen/reopen judgment rather than reopen this note merely to restate a stable answer-side handoff or paired cutover judgment \\
\bottomrule
\end{tabularx}
\caption{A minimal public challenge-stop-profile card. This is the smallest semantic stop object later terse papers can quote directly when the live issue is which piece one question family decomposes into and which approved narrow field may discharge that piece.}
\label{tab:minimal-public-challenge-stop-profile-card}
\end{table}

This card is intentionally non-decorative: it pins the actual worked question-family keys, semantic-piece keys, and approved stop fields without silently re-owning the earlier audience routing layer or the later exact branch, coverage, surface-hook, and terminal-owner layers that neighboring notes already own.\cite{series:workedexample}

\paragraph{A forced public challenge-stop-profile matrix.}
\begin{table}[t]
\centering
\small
\begin{tabularx}{\linewidth}{@{}lX@{}}
\toprule
Field & Forced public answer \\
\midrule
Public-status sentence floor & \nolinkurl{public_status_sentence} still forces the same stop profile only while the same two semantic pieces remain \nolinkurl{continuity_piece} and \nolinkurl{closure_piece}, and those pieces still resolve only to \nolinkurl{continuity_decision} and \nolinkurl{closure_status} at the same owning paths \\
Package-completeness floor & \nolinkurl{package_completeness} still forces the same stop profile only while \nolinkurl{role_rows_piece} and \nolinkurl{closure_token_piece} remain the semantic split and still resolve only to \nolinkurl{role_rows} and \nolinkurl{closure_status} at the same owning paths \\
Compact-diff floor & \nolinkurl{compact_diff_summary} still forces the same stop profile only while \nolinkurl{diff_delta_piece}, \nolinkurl{recompute_piece}, and \nolinkurl{changed_family_piece} remain the semantic split and still resolve only to \nolinkurl{observed_diffs}, \nolinkurl{recomputed}, and \nolinkurl{changed_line_items} at the same owning paths \\
Route/branch boundary floor & Those shortcuts remain honest only while audience routes continue to import the same stop-profile key for the same question family and exact piece-keyed escalation walks remain neighboring branch-map objects rather than being silently rewritten here \\
Staleness trigger floor & This matrix goes stale as soon as a question family changes, a semantic piece is added / removed / renamed, an approved terminal field changes, an owning path changes, or the maintained profile ceases to be the same approved semantic decomposition for that question family \\
Escalation pointer & Stop here only while one maintained question family still forces the same semantic-piece split and approved stop fields; reopen Synthesis~36 when the live issue is audience routing, widen to Synthesis~38 when the exact piece-keyed walk is live, and widen to Synthesis~39--41 when coverage, surface realization, or exact terminal ownership has become the honest moving part \\
\bottomrule
\end{tabularx}
\caption{A forced public challenge-stop-profile matrix. This is the smallest semantic stop object later terse papers can quote directly when the live issue is whether one maintained question family still forces the same semantic-piece split and approved narrow stopping fields rather than merely which profile exists in the abstract.}
\label{tab:forced-public-challenge-stop-profile-matrix}
\end{table}

This matrix is intentionally non-decorative: it pins the actual worked question-family keys, semantic-piece sets, and approved stop fields that still force one maintained stop profile, together with the conditions under which that forcing remains honest or goes stale, without silently re-owning the earlier audience routing layer or the later exact branch / coverage / terminal-owner layers.\cite{series:workedexample}

\paragraph{One tiny challenge-stop-profile drill.}
For the worked auditor public-status token question, the maintained stop profile \nolinkurl{public_status_token} still says that the token itself may stop at \nolinkurl{public_status}, while the continuity and closure challenges may stop only at \nolinkurl{continuity_decision} and \nolinkurl{closure_status}. If the live dispute becomes ``which audience should start where?'' the honest stop is no longer this profile and a later terse paper must reopen Synthesis~36. If the live dispute instead stays semantic but the maintained approved field for \nolinkurl{closure_piece} changes, then the stop-profile matrix has gone stale and this note must be reopened rather than quoted as if the old semantic split still held.\cite{series:workedexample}

\begin{thebibliography}{99}\small

\bibitem{series:downstreamnormalforms}
Anonymous.
\newblock Downstream Normal Forms for Successor Maintenance: Summary, Support, Citation, Quote, and Clause Discipline.
\newblock Series synthesis note (Synthesis~31), 2026. (This archive.)

\bibitem{series:sentencelocks}
Anonymous.
\newblock Sentence Locks and Inline Provenance for Successor Maintenance: Pair-Anchored Reuse, Recheck Guards, and Visible Capsules.
\newblock Series synthesis note (Synthesis~35), 2026. (This archive.)

\bibitem{series:challengeroutes}
Anonymous.
\newblock Challenge Routes and Audience Profiles for Successor Maintenance: Audience-Keyed Start Objects and Escalation Order.
\newblock Series synthesis note (Synthesis~36), 2026. (This archive.)

\bibitem{series:challengecoverage}
Anonymous.
\newblock Challenge Coverage Certificates for Successor Maintenance: Audience Coverage, Surface-Carried Pieces, and Support-Only Hooks.
\newblock Series synthesis note (Synthesis~39), 2026. (This archive.)

\bibitem{series:challengesurfacehooks}
Anonymous.
\newblock Challenge Surface Hooks for Successor Maintenance: Piece-to-Segment Realization and Shared Surface Bindings.
\newblock Series synthesis note (Synthesis~40), 2026. (This archive.)

\bibitem{series:challengeterminalwitnesses}
Anonymous.
\newblock Challenge Terminal Witness Ledgers for Successor Maintenance: Narrow Stop Ownership for Piece Hooks and Support-Only Escalation.
\newblock Series synthesis note (Synthesis~41), 2026. (This archive.)

\bibitem{series:challengebranches}
Anonymous.
\newblock Challenge Branch Maps for Successor Maintenance: Piece-Keyed Escalation from Audience Routes to Approved Terminal Fields.
\newblock Series synthesis note (Synthesis~38), 2026. (This archive.)

\bibitem{series:challengeanswerreviewmenus}
Anonymous.
\newblock Challenge Answer Review Menus for Successor Maintenance: Answer-Side Terminal Citation Normal Form.
\newblock Series synthesis note (Synthesis~54), 2026. (This archive.)

\bibitem{series:publicrequestresponsepacketrefreshresponsepacketclosureverdicts}
Anonymous.
\newblock Public Request Response Packet Refresh Response Packet Closure Verdicts for Source-Only Successor Releases: Stable Request-Side Terminal Forms after Visible Refresh.
\newblock Series synthesis note (Synthesis~73), 2026. (This archive.)

\bibitem{series:pairedterminalcitationdiscipline}
Anonymous.
\newblock Paired Terminal Citation Discipline for Review Spines: Stable Answer Stops and Adjacent Request Widening.
\newblock Series synthesis note (Synthesis~74), 2026. (This archive.)

\bibitem{series:pairedterminalcutoverrules}
Anonymous.
\newblock Paired Terminal Cutover Rules for Review Spines: When Later Terse Papers Stop, Widen, or Reopen.
\newblock Series synthesis note (Synthesis~75), 2026. (This archive.)

\bibitem{series:workedexample}
Anonymous.
\newblock Worked Example for Anonymous-DHT Receipts: Minimal Public Surface, Cross-Release Diff, and Successor Interlock.
\newblock Series synthesis note (Synthesis~17), 2026. (This archive.)

\end{thebibliography}
\end{document}
